% ============================================================
%  Section 8 -- Conclusion
% ============================================================

\chapter{Conclusion}

Le projet \bibopsname{} répond à un besoin concret : évaluer de façon rigoureuse
des LLM et architectures agentiques pour le support IT. Le travail réalisé ne se
limite pas à une démonstration conversationnelle. Il fournit une CLI, un agent
outillé, des outils MCP, des adaptateurs A2A, un moteur d'évaluation
multi-critères, une Racing Arena adversariale, des graphiques et une suite de
tests.

Le résultat principal est favorable à l'architecture agentique : l'agent ReAct atteint \metric{91,34/100} de score composite et
un verdict \verdict{PASS}, contre \metric{46,14/100} et \verdict{FAIL} pour le
LLM unique. Il améliore la qualité moyenne, réduit le coût, la latence, les
tokens et l'empreinte estimée. Ces gains montrent qu'un modèle plus petit peut
devenir plus utile lorsqu'il est placé dans un système outillé et mesurable.

La conclusion reste prudente : les résultats dépendent du dataset, du juge, des
prompts, du corpus RAG et de la configuration expérimentale. Avant toute
industrialisation, il faut augmenter le corpus, introduire une évaluation
humaine, renforcer la sécurité des endpoints et automatiser la régénération des
rapports. La Racing Arena rappelle aussi qu'un agent utile est une surface
d'attaque : les outils doivent être contrôlés, tracés et validés.

Le projet pose donc une base solide pour la suite. Il permet de mesurer,
comparer puis optimiser des LLM pour le support IT, avec une méthode
reproductible et compatible avec les contraintes d'une plateforme de type
BibOps.
